% GENERATED FILE. Edit canonical lessons, then run npm run generate:books.
% canonical-source-hash: fnv1a64:e94ca7143281a453
% canonical-lessons: TE-C100-kudi, TE-C100-edama, TE-C100-tinnaga, TE-C100-daggara, TE-C100-duram

\chapter{Right, Left, Straight ahead, Near, Far}
\label{ch:te-right-left-straight-ahead-near-far}
\hlchaptermodality{100}

\begin{chapteropening}
\textbf{By the end of this chapter:} \emph{I can say right, left, straight ahead, near and far.}

Complete the last lesson of chapter 100: i can say right, left, straight ahead, near and far.
\end{chapteropening}

\section[kuḍi]{\te{కుడి} (kuḍi) — right (not left)}

\label{lesson:TE-C100-kudi}



\begin{quote}
Before the new one: say the Telugu for a gift, a prize, then the Telugu for an invitation.
\end{quote}

\subsection*{You'll want to know: \te{కుడి}}
\textbf{\te{కుడి}} — \emph{kuḍi} — \textquotedblleft{}right (not left)\textquotedblright{}.

\begin{grammarlens}[title={What you've built}]
One more word for this chapter.
\end{grammarlens}

\subsection*{Your turn}
\begin{itemize}
\raggedright
  \item \emph{Say it:} \emph{kuḍi}
  \item \emph{Say it:} \emph{kuḍi}, once more
  \item \emph{Say it:} say \emph{kuḍi}
  \item \emph{Recall:} say the Telugu for a gift, a prize, then the Telugu for an invitation, then say \emph{kuḍi} again
\end{itemize}

\begin{tcolorbox}[breakable,colback=teal!4,colframe=teal!35!black,title={Before you move on}]
What does \te{కుడి} mean? (\textquotedblleft{}Right (not left)\textquotedblright{}.) Say it once more.
\end{tcolorbox}
\section[eḍama]{\te{ఎడమ} (eḍama) — left (not right)}

\label{lesson:TE-C100-edama}



\begin{quote}
Before the new one: say the Telugu for an invitation, then the Telugu for right.
\end{quote}

\subsection*{You'll want to know: \te{ఎడమ}}
\textbf{\te{ఎడమ}} — \emph{eḍama} — \textquotedblleft{}left (not right)\textquotedblright{}.

\begin{grammarlens}[title={What you've built}]
One more word for this chapter.
\end{grammarlens}

\subsection*{Your turn}
\begin{itemize}
\raggedright
  \item \emph{Say it:} \emph{eḍama}
  \item \emph{Say it:} \emph{eḍama}, once more
  \item \emph{Say it:} say \emph{eḍama}
  \item \emph{Recall:} say the Telugu for an invitation, then the Telugu for right, then say \emph{eḍama} again
\end{itemize}

\begin{tcolorbox}[breakable,colback=teal!4,colframe=teal!35!black,title={Before you move on}]
What does \te{ఎడమ} mean? (\textquotedblleft{}Left (not right)\textquotedblright{}.) Say it once more.
\end{tcolorbox}
\section[tinnagā]{\te{తిన్నగా} (tinnagā) — straight ahead}

\label{lesson:TE-C100-tinnaga}



\begin{quote}
Before the new one: say the Telugu for right, then the Telugu for left.
\end{quote}

\subsection*{You'll want to know: \te{తిన్నగా}}
\textbf{\te{తిన్నగా}} — \emph{tinnagā} — \textquotedblleft{}straight ahead\textquotedblright{}.

\begin{grammarlens}[title={What you've built}]
One more word for this chapter.
\end{grammarlens}

\subsection*{Your turn}
\begin{itemize}
\raggedright
  \item \emph{Say it:} \emph{tinnagā}
  \item \emph{Say it:} \emph{tinnagā}, once more
  \item \emph{Say it:} say \emph{tinnagā}
  \item \emph{Recall:} say the Telugu for right, then the Telugu for left, then say \emph{tinnagā} again
\end{itemize}

\begin{tcolorbox}[breakable,colback=teal!4,colframe=teal!35!black,title={Before you move on}]
What does \te{తిన్నగా} mean? (\textquotedblleft{}Straight ahead\textquotedblright{}.) Say it once more.
\end{tcolorbox}
\section[daggara]{\te{దగ్గర} (daggara) — near}

\label{lesson:TE-C100-daggara}



\begin{quote}
Before the new one: say the Telugu for left, then the Telugu for straight ahead.

One slot changes: \textbf{\te{ఇక్కడ}}, \textbf{\te{అక్కడ}}, \textbf{\te{ఎక్కడ}}. Here, there, where?
\end{quote}

\subsection*{You'll want to know: \te{దగ్గర}}
\textbf{\te{దగ్గర}} — \emph{daggara} — \textquotedblleft{}near\textquotedblright{}.

\begin{grammarlens}[title={What you've built}]
One more word for this chapter.
\end{grammarlens}

\subsection*{Your turn}
\begin{itemize}
\raggedright
  \item \emph{Say it:} \emph{daggara}
  \item \emph{Say it:} \emph{daggara}, once more
  \item \emph{Say it:} say \emph{daggara}
  \item \emph{Recall:} say the Telugu for left, then the Telugu for straight ahead, then say \emph{daggara} again
\end{itemize}

\begin{tcolorbox}[breakable,colback=teal!4,colframe=teal!35!black,title={Before you move on}]
What does \te{దగ్గర} mean? (\textquotedblleft{}Near\textquotedblright{}.) Say it once more.
\end{tcolorbox}
\section[dūraṁ]{\te{దూరం} (dūraṁ) — far; distance}

\label{lesson:TE-C100-duram}



\begin{quote}
Before the new one: say the Telugu for straight ahead, then the Telugu for near.
\end{quote}

\subsection*{You'll want to know: \te{దూరం}}
\textbf{\te{దూరం}} — \emph{dūraṁ} — \textquotedblleft{}far; distance\textquotedblright{}.

\begin{grammarlens}[title={What you've built}]
One more word for this chapter.
\end{grammarlens}

\subsection*{Your turn}
\begin{itemize}
\raggedright
  \item \emph{Say it:} \emph{dūraṁ}
  \item \emph{Say it:} \emph{dūraṁ}, once more
  \item \emph{Say it:} say \emph{dūraṁ}
  \item \emph{Recall:} say the Telugu for straight ahead, then the Telugu for near, then say \emph{dūraṁ} again
\end{itemize}

\begin{tcolorbox}[breakable,colback=teal!4,colframe=teal!35!black,title={Before you move on}]
What does \te{దూరం} mean? (\textquotedblleft{}Far; distance\textquotedblright{}.) Say it once more.
\end{tcolorbox}
